\documentclass[a4paper, 11pt]{report}% autres choix : book, report
\usepackage[utf8]{inputenc}% gestion des accents (source)
\usepackage[T1]{fontenc}% gestion des accents (PDF)
\usepackage[francais]{babel}% gestion du français

\usepackage{textcomp}% caractères additionnels
\usepackage{mathtools,amssymb,amsthm}% packages de l'AMS + mathtools
\usepackage{lmodern}% police de caractère
\usepackage{geometry}% gestion des marges
\usepackage{graphicx}% gestion des images
\usepackage{xcolor}% gestion des couleurs
\usepackage{array}% gestion améliorée des tableaux
\usepackage{calc}% syntaxe naurelle pour les calculs
\usepackage{titlesec}% pour les sections
\usepackage{titletoc}% pour la table des matières
\usepackage{fancyhdr}% pour les en-têtes
\usepackage{titling}% pour le titre
\usepackage{enumitem}% pour les listes numérotées
\usepackage[hyphens]{url} % Pour des césures correctes dansles URLs
\usepackage[pdfauthor = {{MB TC AP}}, pdftitle = {{Modelisation Lokta Volterra}}, pdfstartview = Fit, pdfpagelayout = SinglePage , pdfnewwindow = true, bookmarksnumbered = true, breaklinks, colorlinks, linkcolor = black, urlcolor = black, citecolor = cyan, linktoc = all]{hyperref}
\usepackage{titlesec}

\titleformat{\chapter}[block]
  {\normalfont\Huge\bfseries}% font of number
  {\thechapter~ -}% format of number
  {20pt}% space between number and title
  {\Huge}% font of title
  
% Taille des marges
\geometry{hmargin=2.8cm,vmargin=2.8cm}

% RAZ des numéros de section après un chapitre
\makeatletter\@addtoreset{section}{chapter}\makeatother
% Pour mettre des I, II, etc. aux parties
\renewcommand{\thepart}{\Roman{part}}
% Pour mettre des 1, 2, etc. aux chapitres
\renewcommand{\thechapter}{\Roman{chapter}}
% Idem pour les sections et avoir le numéro de chapitre
\renewcommand{\thesection}{\Roman{section}}
% Idem pour les sous-sections et avoir le numéro de chapitre
\renewcommand{\thesubsection}{\Roman{section}.\arabic{subsection}}

\titlecontents{section}%
	[1.5em]% retrait à gauche
	{\addvspace{1em plus 0pt}\bfseries}% matériel avantcommun aux entrées numérotées ou pas
	{\contentslabel{2em}}% avant lorsqu'il y a un numéro
	{\hspace{-1.3em}}% avant lorsqu'il n'y a pas de numéro
	{\hfill\contentspage}% points de suspension et no page
	[\addvspace{0pt}]% matériel après
	
	\dottedcontents{subsection}%
	[3em]% retrait gauche
	{\addvspace{0pt}}% matériel avant
	{1.5em}% espacement de contentslabel
	{0.75em}% espace entre les . . . .
    [\addvspace{0pt}]% matériel après

    \dottedcontents{subsubsection}%
	[4.5em]% retrait gauche
	{\addvspace{0pt}}% matériel avant
	{1.5em}% espacement de contentslabel
	{0.75em}% espace entre les . . . .
    [\addvspace{0pt}]% matériel après

% En-tête, pied de page
\pagestyle{fancy}
\fancyhead[O]{Mini-Pojet \\Lokta Volterra}
\fancyhead[R]{Modélisation ST2}
\fancyhead[L]{\includegraphics[height=1cm]{images/cs.png}}

\renewcommand{\footrulewidth}{1pt}
\fancyfoot[C]{\thepage}
\fancyfoot[R]{Matthieu BRIET \\Tanguy COLLEVILLE \\ Antoine PAGNEUX}
\fancyfoot[L]{Groupe 218}
    

\begin{document}
    % Page de Garde
    \begin{titlepage}
        \newcommand{\HRule}{\rule{\linewidth}{0.5mm}}
        \begin{center}
            % En-têtes
            \textsc{\LARGE{} Mini-Projet Modélisation} \\[0.5cm] 
            \textsc{\Large{} CentraleSupélec - 1A} \\[0.5cm]
            \textsc{\large{} Groupe 218} \\[0.5cm] 
            % Titre
            \HRule \\[0.6cm]
            {\huge\bfseries{} Lokta Volterra} \\[0.25cm]
            \HRule \\[1.5cm]
            % Date
            {\large\today} \\[2cm] 
            % Logo CS
            \includegraphics[width=8cm]{images/cs.png}
            \\[2cm] 
        \end{center}
        \vfill{}
        % Auteurs
        \begin{minipage}{0.45\linewidth}
            \begin{flushleft}
                \Large\textit{Auteurs :} \\
                Matthieu \textsc{BRIET} \\
                Tanguy \textsc{COLLEVILLE} \\
                Antoine \textsc{PAGNEUX} 
            \end{flushleft}
        \end{minipage}
    \end{titlepage}

    % Table des matières
    \tableofcontents
    % Table des figures
    \listoffigures
    % Table des tableaux
    %\listoftables

    \newpage
    \section{Introduction}
        \subsection{blabla}
        \subsection{blabla}

    \section{Modélisation à état continu}
        \subsection{blabla}

    \section{Modélisation à événements discrets}
        \subsection{blabla}


\end{document}